\section{Introduction}

A critical question in understanding the labor market and informing policy recommendations is whether the observed differences in earnings between formal and informal employment results from market segmentation or incentives within a competitive market framework. In other words, do individuals opt for informal employment by choice, or are they pushed into it due to barriers to entry in the formal employment? This question is particularly important in developing economies, where informality is widespread, yet no consensus has been achieved on the driving mechanisms to understand the complex structure of informal employment.

Traditional dual market theory, beginning with Lewis (1954), conceives informal employment as a backward, marginal, and disadvantaged sector where employees are paid wages above market-clearing prices and use it as a last resort to escape unemployment (Harris and Todaro, 1970; Fields, 1990, 2005a; Stiglitz, 1976). According to this approach, workers in the informal sector are predominantly involuntarily engaged and earn less than identical workers in the formal sector. This segmentation may result from access to information, geographic barriers, or sector-specific variations in employment and wage regulations (Dickens and Lang, 1992). For instance, firms may offer wages above market-clearing prices to increase retention and reduce turnover costs, ultimately maximizing employee productivity, raising the price of labor, and reducing demand for workers in the formal sector.

The notion of informal employment as a transitional ground for entry into formal employment was challenged by the International Labour Organization (ILO, 1973), which viewed the informal sector as a permanent, economically sufficient, and profit-making small-scale urban economic continuum of employment and income, with little evidence of sectoral mobility (Sethuraman, 1981). This perspective gave rise to an alternative theory suggesting that informal employment results from workers' voluntary choices within a competitive market framework (Maloney, 1999; Rosenzweig, 1988). Working in the informal sector reflects a worker's choice based on the cost-benefit analysis of each sector, given their individual characteristics. Voluntary employment in the informal sector can be due to workers trading pecuniary rewards for non-pecuniary benefits, maximizing utility rather than earnings (Maloney, 2004), self-selecting into jobs based on comparative advantage (Roy, 1951), or seeking to avoid tax burdens and inefficient regulations (Fields, 2003).

Thus, two theories explain informal employment: the segmentation hypothesis views it as a survival strategy for those excluded from formal jobs, while the comparative advantage hypothesis sees it as a voluntary choice by workers to maximize their utility or earnings. While the differential between formal and informal employment is well-documented (Mazumdar, 1975; Pradhan and Van Soest, 1995; Badaoui et al., 2008; Blunch, 2015), lower wages and returns to education and experience in informal employment do not necessarily imply market segmentation as long as there is mobility between these two sectors (Gunther and Launov, 2012). The observed differences in wage distributions may be associated with non-random selection processes or penalties due to differences in workers' skills or abilities. Therefore, it is necessary to account for the unobservable characteristics of workers that determine sector choices and earnings to distinguish between segmentation and comparative advantage. A range of empirical studies has explored whether labor markets in developing and developed countries operate competitively or are segmented (Dickens and Lang, 1985; Heckman and Hotz, 1986; Rosenzweig, 1988; Magnac, 1991; Pratap and Quintin, 2006).

The empirical literature on informal employment predominantly assumes the informal sector is homogeneous (Gunther and Launov, 2012). However, theoretical studies have proposed a heterogeneous informal sector with two distinct tiers: the lower-tier and upper-tier (Fields, 1990). The lower-tier fits the dual market model, assuming formal employment pays higher wages and provides job security and social security contributions, whereas informal employment involves individuals who cannot afford unemployment, are paid low wages, and are excluded from formal jobs. The upper-tier comprises individuals who voluntarily choose informal employment, expected to earn more than formal employment due to their specific characteristics and non-job benefits (Cunningham and Maloney, 2001; Henley et al., 2009).

The level of wages and earnings of the workforce is determined by the forces of labor demand and supply based on the neo-classical framework. However, this is not true for developing countries, where labor market dualism and strong entry barriers exist across different segments of the labor market. The dual labor market model posits the existence of two distinct sectors of economic activity, usually classified as formal and informal, extending to the nature of employment. Within formal employment, employees have stable jobs with higher pay, better working conditions, permanent contracts, paid leave, and social security contributions provided by the employer. Although informal employment lacks employment contracts and job-based benefits, it offers more flexibility in terms of working hours and locations. This asymmetry affects how employees value formal and informal employment differently.

Those in poverty are more likely to participate in informal employment, where wage rates and total income are lower than formal employment, further reinforcing their marginalized position (Williams, 2014). The productivity level of informal employment is lower than that of formal employment due to differences in technology, lack of transparency, reliance on family sources for credit, and limitations in acquiring modern management skills and worker training (Monsted, 2020; Benjamin et al., 2014). These factors limit growth potential, access to the global market, and result in fragile working conditions. Because of these productivity mismatches and different wage-setting mechanisms, the extent of wage discrimination between formal and informal employees varies based on personal, household, job, and employer characteristics. Formal employment is governed by anti-discrimination and equal treatment laws, whereas these regulations are absent in informal employment, with competition and social norms favoring dominant groups (Neog and Sahoo, 2016). Understanding wage discrimination within the context of distinct employment opportunities and characteristics is crucial to comprehending labor market dynamics and welfare distribution in the economy.

However, few studies exist that have analyzed this theory of heterogeneous informal market along with their distribution from an empirical perspective. The existing literature often focuses on gender wage gaps, gender gaps within informality,  but lacks studies on wage gaps related to job occupations in informality. As a result, we cannot directly observe, whether employees within distinct job occupations view informal employment as a opportunity or as a strategy of the last resort (Gunther, Launov, 2012). The limitations of our study is that, we only look at the wage employees and exclude self-employment, which provides a narrow selection of possible informal employment.

The main objective of our paper is to examine the impact of informality on earnings for employees in formal and informal sector after considering the unobserved characteristics. To this purpose, we use two rounds of household-level survey data \footnote{Nepal Labor Force Survey round II (2008) and round III (2018)}, collected by Nepal Statistics Office, and compare the earnings differences between formal and informal employment workers within three tiers of occupations across the distribution. This paper focuses on individuals aged 15-60 years who are in full-time wage employment, defined as those working in private financial and non-financial institutions, non-profit institutions, and other entities for 40 hours or more per week with cash as the mode of payment. We exclude agriculture sector in our study, We aim to examine whether significant gaps in earnings exist between formal and informal wage earners. The choice of full-time wage employees is motivated by several reasons: it is important to analyze employees separately from other groups, such as self-employed workers, because determinants of earnings differences, such as employer and job characteristics, apply only to full-time wage employees; the definition of informality varies across employment statuses \footnote{We follow the conceptual framework of informal employment defined by the 15th International Conference of Labour Statisticians as used in the Nepal Labor Force Survey III }. Based on these,  we analyze if there is heterogeneity in the informal sector and how this is related to earnings and distinct job sectors. We identify three groups of job sector workers: elementary, clerical and managers respectively. In order to find the wage differences across the distribution, this paper uses linear Recentered Influence Function (RIF) - OLS regressions to estimate the Mincer-type wage equations for distinct occupations (Firpo, 2009;2018). Estimation at each quantiles is based on the Unconditional Quantile Regressions (UQR) . This method allows us to compute a detailed decomposition in a path-independent way and allows for the unconditional mean interpretation of the coefficient estimates, in contrary to the conditional quantile regressions. For robustness purpose, in the concerns of nonlinearity, we apply the reweighting scheme proposed by \textcite{DiNardo1996} and show that the results do not change significantly.  


Our findings indicate a dual market structure in the Nepalese labor market, revealing a consistent wage penalty for informal sector workers across various occupations in 2008. There exists a heterogeneous Specifically, informal employees in clerical, elementary, and managerial roles earn less than their formal counterparts, even after controlling for nonlinearity in the model. In the elementary occupation group, a wage premium for informal employment is observed up to the 40th quantile, likely due to the scarcity of formal jobs in these lower quantiles. Conversely, clerical and managerial roles exhibit a more pronounced wage penalty for informal workers, particularly beyond the 40th quantile and 20th quantile respectively. Comparing 2008 and 2018, the informal wage penalty diminishes, suggesting wage convergence over time. For clerical workers, the wage gap is most significant around the median in 2008 but decreases notably by 2018. In the managerial group, the wage disparity peaks at the 20th quantile in 2008 and decreases steadily across the distribution, with the most significant gap between the 20th and 40th quantiles in 2018, showing signs of convergence. Overall, while there is a persistent wage penalty for informal workers, this gap is narrowing over time, highlighting potential improvements in labor market dynamics. The literature shows the phenomenon of glass ceiling and sticky floor for the quantile-specific wage gaps, but there exists no such phenomenon, but wage gap is more pronounced around the median in our study. Similarly, the respective categories contribute differently to the explained(endowment effects) and unexplained(coefficient effects) of the respective pay gaps.

The remainder of this paper is organized as follows: Section 2 presents the development of our econometric model. Section 3 describes the dataset and provides descriptive statistics on wage inequality among formal and informal employees in Nepal. Section 4 discusses the results, focusing on selection into OLS-RIF regression and decomposition based on quantile regression for formal and informal employment in Nepal. Finally, Section 5 concludes the paper with potential avenues for further research.
